\documentclass[11pt]{article}
\usepackage[margin=0.8in]{geometry}
\usepackage[T1]{fontenc}
\usepackage{lmodern,amsmath,amssymb,amsthm,booktabs,longtable}
\usepackage[colorlinks=true,linkcolor=blue,urlcolor=blue,citecolor=blue]{hyperref}
\hypersetup{pdftitle={Arithmetic identities, decimal periods, and complementary Fibonacci limits},pdfauthor={Blake Stone},pdfsubject={Collected OEIS proofs with Lean verification}}
\setlength{\parindent}{0pt}
\setlength{\parskip}{0.4em}
\newtheorem{theorem}{Theorem}[section]
\newtheorem{lemma}[theorem]{Lemma}
\newcommand{\N}{\mathbb N}
\newcommand{\R}{\mathbb R}
\begin{document}
\begin{center}
{\LARGE\bfseries Arithmetic identities, decimal periods,\\[3pt]
and complementary Fibonacci limits}\par
\vspace{0.5em}
{\large Blake Stone}\par
Proof developed with OpenAI Astra 6.2\par
November 26, 2026
\end{center}

\textbf{Abstract.}
We prove a unitary-prime-divisor identity, a squared-valuation formula for a
weighted binary digit sum, a conjectured decimal period, and an
integer-square property for generalized binomial coefficients. For the binary
digit sum we also refute a proposed range for its asymptotic constant under
natural logarithms. A further uniform argument proves golden-ratio
convergence for the 34 complementary-sequence equations A294532--A294565,
with an explicit error bound. The accompanying Lean 4.19.0 packages check
the stated identities, limits, and disproof. The family theorem quantifies
over sequences satisfying the exact defining complementary equations.

\section{A143132: a four-term decimal period}

OEIS A143132 is the binomial transform of $1,5,15,35,70$, followed by zeros,
with its first term indexed by $1$~\cite{A143132}. Its comment conjectures
the repeating last-digit pattern $1,6,6,6$.

\begin{theorem}
For every integer $n\ge1$,
\[
 a(n)\bmod10=
 \begin{cases}1,&n\equiv1\pmod4,\\6,&\text{otherwise}.\end{cases}
\]
\end{theorem}
\begin{proof}
Write $b(m)=a(m+1)$. The defining binomial transform is
\[
 b(m)=1+5m+15\binom m2+35\binom m3+70\binom m4.
\]
Pascal's identity, applied four times, gives the exact integer identity
\begin{equation}
 b(m+4)=b(m)+10\left(32+55m+56\binom m2+28\binom m3\right).
 \label{eq:decimal}
\end{equation}
For example, the coefficients of $1,m,\binom m2,\binom m3$ in the
difference are $320,550,560,280$, respectively; the $\binom m4$ terms
cancel. The identity includes $m<4$, with the convention $\binom mk=0$
when $k>m$.

Thus $b(m+4)\equiv b(m)\pmod{10}$. The initial values are
$b(0),b(1),b(2),b(3)=1,6,26,96$, whose residues are $1,6,6,6$.
Induction and the shift $m=n-1$ prove the theorem.
\end{proof}

The entry already contains generating-function and polynomial descriptions
by R. J. Mathar and T. D. Noe from 2008. Those formulas are prior work.
The result here is the explicit proof and formal verification of the
listed congruence assertion, not a claim that the polynomial is new.

\section{A063956: the unitary-prime product identity}

Let $u(n)$ be the sum of the primes whose exponent in the factorization of
$n$ is exactly one. These are precisely its unitary prime divisors.
Velin Yanev's February 17, 2019 conjecture in A063956~\cite{A063956} is
the following identity.

\begin{theorem}
For all positive integers $n,m$,
\begin{equation}
 u(nm)=u(n)+u(m)-u(\gcd(n^2,m))-u(\gcd(n,m^2)).
 \label{eq:unitary}
\end{equation}
\end{theorem}
\begin{proof}
Fix a prime $p$, and put $x=\nu_p(n)$, $y=\nu_p(m)$. The respective
exponents in $nm$, $\gcd(n^2,m)$, and $\gcd(n,m^2)$ are
$x+y$, $\min(2x,y)$, and $\min(x,2y)$.

With $[P]$ denoting $1$ when $P$ is true and $0$ otherwise, we have
\begin{equation}
 [x+y=1]+[\min(2x,y)=1]+[\min(x,2y)=1]=[x=1]+[y=1].
 \label{eq:local}
\end{equation}
If one of $x,y$ is zero, the two minimum terms vanish, and the identity
is immediate. If both are positive, the first term vanishes,
$\min(2x,y)=1$ is equivalent to $y=1$, and $\min(x,2y)=1$ is equivalent
to $x=1$.

Multiply \eqref{eq:local} by $p$ and sum over the finitely many primes
dividing $nm$. The result is
\[
 u(nm)+u(\gcd(n^2,m))+u(\gcd(n,m^2))=u(n)+u(m),
\]
which rearranges to \eqref{eq:unitary}. It also proves that the two
subtractions on its right are legitimate over the natural numbers.
\end{proof}

The Lean proof uses the standard prime-factorization function, proves the
local exponent identity, sums it over the union of the finite supports,
and substitutes the standard multiplication, squaring, and gcd rules.

\section{A073642: binary carries and squared valuations}

For $n=\sum_{j\ge0}\varepsilon_j2^j$, with $\varepsilon_j\in\{0,1\}$,
write
\[
 s_2(n)=\sum_j\varepsilon_j,\qquad w(n)=\sum_jj\varepsilon_j.
\]
The function $w$ is A073642. Yanev's September 9, 2017 conjecture in that
entry~\cite{A073642} relates it to the squares of the binary valuations.

\begin{theorem}
For every integer $n\ge0$,
\begin{equation}
 \boxed{2w(n)+\sum_{k=1}^n\nu_2(k)^2+3s_2(n)=3n.}
 \label{eq:binary}
\end{equation}
Equivalently,
\[
 w(n)=\frac{3(n-s_2(n))-\sum_{k=1}^n\nu_2(k)^2}{2}.
\]
\end{theorem}
\begin{proof}
Fix $k\ge1$, and set $t=\nu_2(k)$. Incrementing $k-1$ to $k$ changes
the $t$ trailing one-bits in positions $0,\ldots,t-1$ into zeros, and
changes the bit in position $t$ into one. The higher bits are unchanged.
Consequently,
\begin{align}
 s_2(k)-s_2(k-1)&=1-t,\label{eq:ones-increment}\\
 w(k)-w(k-1)&=t-\sum_{j=0}^{t-1}j
             =\frac{3t-t^2}{2}.\label{eq:weight-increment}
\end{align}
Summing \eqref{eq:ones-increment} from $1$ through $n$ yields
$\sum_{k=1}^n\nu_2(k)=n-s_2(n)$.
Summing twice \eqref{eq:weight-increment} and substituting this identity
gives \eqref{eq:binary}. In particular, the numerator of the divided
formula is nonnegative and even.
\end{proof}

The formalization proves the binary recursions and the two increment
identities, then telescopes. It also proves that its recursive descriptions
agree with the standard \texttt{Nat.digits 2 n} and
\texttt{padicValNat 2 n}. Thus the final theorem concerns the specified
binary digit sum and the ordinary exponent of $2$, rather than auxiliary
functions with an unproved correspondence.

\subsection{A separate asymptotic claim is false}

The same entry suggests that
\[
 S(N):=\sum_{n=1}^Nw(n)\sim C N(\log N)^2,
 \qquad 0.47<C<0.5.
\]
Here and below, $\log$ means the natural logarithm. We refute the proposed
range by proving a subsequence limit; we do not need to assume or prove
existence of the full normalized limit.

Among the $2^m$ binary strings of length $m$, every position is one in
exactly $2^{m-1}$ strings. Hence, for $m\ge1$,
\[
 \sum_{n=0}^{2^m-1} w(n)
 =2^{m-1}\sum_{j=0}^{m-1}j
 =2^{m-2}m(m-1).
\]
Also $w(2^m)=m$, so
\begin{equation}
 S(2^m)=2^{m-2}m(m-1)+m.
 \label{eq:block}
\end{equation}
Since $\log(2^m)=m\log2$, division gives
\[
 \frac{S(2^m)}{2^m(\log(2^m))^2}
 =\frac{1-1/m}{4(\log2)^2}
   +\frac{2^{-m}}{m(\log2)^2}
 \longrightarrow\frac{1}{4(\log2)^2}.
\]
This constant exceeds $1/2$. For an exact bound, the positive-term
exponential series gives
\[
 e^{7/10}\ge 1+\frac7{10}+\frac{(7/10)^2}{2}
              +\frac{(7/10)^3}{6}>2,
\]
so $0<\log2<7/10$ and $1/(4(\log2)^2)>25/49>1/2$.
Any full normalized limit would have to equal this subsequence limit.
Therefore no $C<1/2$ can satisfy the proposed asymptotic.

Lean checks \eqref{eq:block} through exact doubled-range sums, the
subsequence limit, the strict lower bound on its constant, and the
exclusion of every proposed global limit below $1/2$. The scope of this
additional result is precisely a subsequence limit and a disproof of the
stated range, not a full all-$N$ asymptotic theorem.

\section{A294532--A294565: a family of golden-ratio limits}

Let $(a_n)$ and $(b_n)$ be strictly increasing sequences of nonnegative
integers whose ranges are disjoint and together comprise all positive
integers. Their initial values are
\[
 a_0=1,\qquad a_1=2,\qquad b_0=3.
\]
Consider the equations
\begin{equation}
 a_{n+2}=a_{n+1}+a_n+e_n,\qquad
 e_n=p b_{n+1}+q b_n+r n+s\quad(n\ge0),
 \label{eq:family}
\end{equation}
with the coefficients in Table~\ref{tab:family}. These are the 34
individual title equations in OEIS A294532--A294565, rewritten with the
index shifted by two~\cite{A294532}. Each entry conjectures that the ratio
of consecutive $a$ terms tends to the golden ratio.

\begin{theorem}\label{thm:family}
For every pair of sequences satisfying these defining conditions and any
coefficient row in Table~\ref{tab:family}, let $\phi=(1+\sqrt5)/2$. Then
\begin{equation}
 \boxed{\displaystyle \frac{a_{n+1}}{a_n}\longrightarrow\phi.}
 \label{eq:ratio}
\end{equation}
More precisely, for every $n\ge0$,
\begin{equation}
 \left|\frac{a_{n+1}}{a_n}-\phi\right|
 \le97(n+1)\left(\frac23\right)^n.
 \label{eq:error}
\end{equation}
\end{theorem}

\begin{center}
\small
\begin{longtable}{lrrrr@{\qquad}lrrrr}
\caption{Coefficients in the shifted equation \eqref{eq:family}.}\label{tab:family}\\
\toprule
Entry&$p$&$q$&$r$&$s$&Entry&$p$&$q$&$r$&$s$\\
\midrule
\endfirsthead
\toprule
Entry&$p$&$q$&$r$&$s$&Entry&$p$&$q$&$r$&$s$\\
\midrule
\endhead
A294532&0&1&0&0&A294549&1&0&1&3\\
A294533&0&1&0&1&A294550&1&1&0&0\\
A294534&0&1&0&2&A294551&1&1&0&1\\
A294535&0&1&0&3&A294552&1&1&1&2\\
A294536&0&1&0&-1&A294553&1&1&-1&-2\\
A294537&0&1&1&2&A294554&1&1&0&2\\
A294538&0&1&2&4&A294555&1&1&0&3\\
A294539&0&1&1&1&A294556&1&1&1&3\\
A294540&0&1&2&3&A294557&1&1&1&1\\
A294541&1&0&0&0&A294558&1&1&2&4\\
A294542&1&0&0&1&A294559&1&2&0&0\\
A294543&1&0&0&2&A294560&2&2&0&0\\
A294544&1&0&0&3&A294561&2&1&0&0\\
A294545&1&0&0&-1&A294562&1&-1&0&1\\
A294546&1&0&1&2&A294563&1&-1&1&2\\
A294547&1&0&2&4&A294564&2&-1&0&-1\\
A294548&1&0&1&1&A294565&2&-1&0&-2\\
\bottomrule
\end{longtable}
\end{center}

The guide in A294532 lists the final A294565 equation with constant
$-3$, whereas A294565's own title uses $-2$. The proof covers both:
the additional guide variant has $(p,q,r,s)=(2,-1,0,-3)$.
It is recorded as an alternative specification, not counted as another
OEIS sequence.

\subsection{The complement has a linear upper bound}

Strict increase and $b_0=3$ imply $b_n\ge n+3$ and
$b_{n+1}\ge b_n+1$. Substitution into each coefficient row shows
\begin{equation}
 e_n\ge2.\label{eq:force-positive}
\end{equation}
For the rows involving subtraction, use these two inequalities before
simplifying. For example,
$b_{n+1}-b_n+1\ge2$, and
$2b_{n+1}-b_n-3\ge b_n-1\ge n+2$.
Thus \eqref{eq:family} implies $a_{n+2}\ge a_{n+1}+2$.
Apart from the initial pair $1,2$, no two consecutive integers can both
belong to the $a$ sequence.

If $b_{n+1}>b_n+2$, neither $b_n+1$ nor $b_n+2$ is a $b$ term, since
they lie strictly between consecutive $b$ terms. Complementarity places
both in the $a$ sequence. They are at least $4$, contradicting the
preceding gap property. Hence $b_{n+1}\le b_n+2$, and induction gives
\begin{equation}
 n+3\le b_n\le2n+3.\label{eq:b-bounds}
\end{equation}
Substitute these bounds into the finitely many signed expressions in
Table~\ref{tab:family}. They give the uniform bound
\begin{equation}
 0\le e_n\le32(n+1).\label{eq:linear-force}
\end{equation}
Both coefficient bounds are proved for every row in Lean.

\subsection{A contracting error proves convergence}

We give the general analytic argument, which applies to any recurrence
$a_{n+2}=a_{n+1}+a_n+e_n$ with $a_0\ge1$, $a_1\ge3/2$, and
$0\le e_n\le C(n+1)$, where $C\ge0$.

First, a two-step induction gives
\begin{equation}
 a_n\ge(3/2)^n.\label{eq:exponential}
\end{equation}
The induction works because $(3/2)^2\le3/2+1$ and $e_n\ge0$.
In particular, every denominator used below is positive.

Set $d_n=a_{n+1}-\phi a_n$. Since $\phi^2=\phi+1$, we have
\[
 d_{n+1}=(1-\phi)d_n+e_n.
\]
The elementary bounds $3/2<\phi<5/3$ imply $|1-\phi|\le2/3$.
Put $K=3C+|d_0|$. Induction gives
\begin{equation}
 |d_n|\le K(n+1),\label{eq:step-error}
\end{equation}
because the induction step is bounded by
\[
 \frac23 K(n+1)+C(n+1)\le K(n+1)\le K(n+2).
\]
Combining \eqref{eq:exponential} and \eqref{eq:step-error},
\[
 \left|\frac{a_{n+1}}{a_n}-\phi\right|
 \le\frac{K(n+1)}{(3/2)^n}\longrightarrow0.
\]
This proves convergence itself; existence of the ratio limit is never an
assumption. In our family, $C=32$, $d_0=2-\phi$, and
$K=96+2-\phi<97$. Therefore \eqref{eq:error} follows, proving
Theorem~\ref{thm:family}.

\subsection{Relation to the 2019 theorem and formal scope}

Kimberling and Moses~\cite{KimberlingMoses} study this broader subject and
list A294532 among their guides. Their Theorem 2.1 identifies a dominant-root
limit under several hypotheses; hypothesis (ii) explicitly assumes that
the ratio limit exists. The argument above supplies a direct proof of
existence for the present family, together with a uniform quantitative
bound. We do not present their conditional identification as a new result.

The Lean family theorem quantifies over arbitrary sequences satisfying the
OEIS defining conditions: the seed values, strict increase, disjoint
positive cover, and the selected equation. No convergence or growth
condition is added to that specification. The formalization proves the
limit for every such solution; it does not separately construct the
sequences by a mex algorithm or prove their existence and uniqueness.
There is one named limit theorem for each of the 34 entries, plus the
guide's alternative final equation.

\newpage
\section{A129453: a generalized binomial quotient is an integer square}

Define the positive integer
\[
 G(n)=\prod_{i=1}^n\prod_{j=1}^n\gcd(i,j),\qquad G(0)=1,
\]
and, for $0\le k\le n$, the rational number
\[
 T(n,k)=\frac{G(n)}{G(k)G(n-k)}.
\]
OEIS A129453 conjectures that $T(n,k)$ is divisible by the ordinary
binomial coefficient and that the quotient is a perfect
square~\cite{A129453}.

\begin{theorem}
For all integers $0\le k\le n$, there is a positive integer $t$ such that
\begin{equation}
 G(n)=\binom nk G(k)G(n-k)t^2.
 \label{eq:square}
\end{equation}
Consequently $T(n,k)=\binom nk t^2$, and in particular
$T(n,k)/\binom nk$ is an integer square.
\end{theorem}
\begin{proof}
Fix a prime $p$. Counting ordered pairs of multiples of $p^e$ gives
\begin{equation}
 \nu_p(G(n))=\sum_{e\ge1}\left\lfloor\frac{n}{p^e}\right\rfloor^2.
 \label{eq:gcd-valuation}
\end{equation}
Indeed, each pair $(i,j)$ contributes once for every $e\ge1$ for which
$p^e$ divides both $i$ and $j$; there are exactly
$\lfloor n/p^e\rfloor^2$ such pairs. All sums below have finite support.
Legendre's factorial formula also gives
\[
 \nu_p\binom nk=\sum_{e\ge1}
 \left(\left\lfloor\frac n{p^e}\right\rfloor
 -\left\lfloor\frac k{p^e}\right\rfloor
 -\left\lfloor\frac{n-k}{p^e}\right\rfloor\right).
\]

For each $e$, write
\[
 A_e=\left\lfloor\frac k{p^e}\right\rfloor,\quad
 B_e=\left\lfloor\frac{n-k}{p^e}\right\rfloor,\quad
 c_e=\left\lfloor\frac n{p^e}\right\rfloor-A_e-B_e.
\]
Since $n=k+(n-k)$, the floor addition rule gives $c_e\in\{0,1\}$.
The contribution to the difference between the exponent in $G(n)$ and
that in $\binom nk G(k)G(n-k)$ is
\[
 (A_e+B_e+c_e)^2-A_e^2-B_e^2-c_e
 =2\bigl(A_eB_e+c_e(A_e+B_e)\bigr),
\]
which is a nonnegative even integer. Put
\[
 E_p=\sum_{e\ge1}\bigl(A_eB_e+c_e(A_e+B_e)\bigr),\qquad
 t=\prod_{\substack{p\le n\\p\text{ prime}}}p^{E_p}.
\]
Every prime therefore has the same exponent on the two sides of
\eqref{eq:square}. All factors are positive, and equality follows from
unique factorization. The boundary cases $n=0$, $k=0$, and $k=n$ have
$t=1$ and are included.
\end{proof}

The exact division follows from \eqref{eq:square} and the positivity of
$G(k)$, $G(n-k)$, and $\binom nk$. The formal proof starts with the
original double gcd product; it does not assume a conjectured
factorization formula as a hypothesis.

\subsection{Attribution and the stronger target}

Ordinary integrality of $T(n,k)$ already follows from Wanli Ma's
2025 preprint, Theorem 2.5 and its square-exponent gcd-product
example~\cite{Ma}. Formula \eqref{eq:gcd-valuation} also appears there,
with credit to the earlier OEIS verification by Charles R. Greathouse IV.
The related identity $G(n)=n!H(n)^2$, where
$H(n)=\prod_{1\le i<j\le n}\gcd(i,j)$, is recorded in A224479 with
Daniel Forgues's 2013 explanation~\cite{A224479}.
These identities and the ordinary-integrality result are prior work.
Our target here is the stronger stated integer-square quotient, together
with a formal proof from the original product definition. Targeted
searches did not locate a prior proof of that precise stronger assertion;
this is not an exhaustive priority claim.

\section{Verification, sources, and reproduction}

The accompanying packages use Lean 4.19.0 and the following Mathlib commit:
\begin{center}
\texttt{c44e0c8ee63ca166450922a373c7409c5d26b00b}.
\end{center}
Each includes a pinned toolchain and dependency manifest, successful fresh
build output, and a separate mathematical and formal review. The final
theorems have only the standard logical dependencies
\texttt{propext}, \texttt{Classical.choice}, and \texttt{Quot.sound}.
There are no unfinished proofs, added axioms, unsafe replacements, or
\texttt{native\_decide} calls.

The principal checked results are:
\begin{itemize}
\item \texttt{A143132.conjecture}: the last-digit pattern with the correct
positive-index offset.
\item \texttt{A063956.conjecture}: the product identity for all positive
arguments.
\item \texttt{A073642.oeis\_conjecture}: the squared-valuation formula,
expressed with standard binary digits and binary valuation.
\item \texttt{A073642.power\_two\_limit} and
\texttt{A073642.no\_limit\_below\_half}: the subsequence limit and the
natural-log asymptotic disproof.
\item \texttt{ComplementaryFibonacci.solution\_ratio\_limit} and
\texttt{solution\_ratio\_error\_bound}: the complete family argument
and its quantitative estimate.
\item \texttt{A129453.square\_quotient} and \texttt{A129453.quotient\_square}: the original
double gcd product equals the binomial coefficient times the two
smaller products times an integer square.
\end{itemize}

The README in each package gives its Mathlib cache command and
the command \texttt{lake build}. The audit records distinguish the exact
formal statements from their explanatory corollaries. The independent
reviews also check sequence definitions, offsets, hypotheses, and prior
work, beyond successful compilation.

The individual family equations were checked against the official
\texttt{oeis/oeisdata} export dated November 26, 2026, including the
A294565 discrepancy. The retrieved entries still label the relevant
assertions as conjectural or tentative. Searches for identifiers and
equivalent formulations did not locate earlier proofs of the targeted
assertions beyond the related results credited above. These are bounded
source checks, not an exhaustive priority claim or a claim of OEIS
review or acceptance.

\newpage
\begin{thebibliography}{9}
\bibitem{A143132} G. W. Adamson, OEIS A143132, July 27, 2008;
formulas by R. J. Mathar and T. D. Noe, 2008.
\url{https://oeis.org/A143132}.
\bibitem{A063956} E. Labos, OEIS A063956, September 4, 2001;
conjecture by V. Yanev, February 17, 2019.
\url{https://oeis.org/A063956}.
\bibitem{A073642} B. Cloitre, OEIS A073642, August 29, 2002;
valuation identity conjectured by V. Yanev, September 9, 2017.
\url{https://oeis.org/A073642}.
\bibitem{A294532} C. Kimberling, OEIS A294532--A294565, November 2017.
Family guide: \url{https://oeis.org/A294532}; final individual equation:
\url{https://oeis.org/A294565}.
\bibitem{KimberlingMoses} C. Kimberling and P. J. C. Moses,
``Linear Complementary Equations and Systems,''
\emph{The Fibonacci Quarterly} 57(5) (2019), 96--110.
\url{https://www.fq.math.ca/Papers1/57-5/kimberling-moses.pdf}.
\bibitem{OEISData} The OEIS Foundation, official Git-formatted export,
snapshot time November 26, 2026.
\url{https://github.com/oeis/oeisdata}.
\bibitem{A129453} P. Bala, OEIS A129453, April 16, 2007.
\url{https://oeis.org/A129453}.
\bibitem{A224479} P. Luschny, OEIS A224479, April 7, 2013;
matrix factorization explanation by D. Forgues, April 14, 2013.
\url{https://oeis.org/A224479}.
\bibitem{Ma} W. Ma, \emph{Generalized factorials characterized by
Dirichlet convolution}, arXiv:2509.13354v1, September 14, 2025.
\url{https://arxiv.org/abs/2509.13354}.
\end{thebibliography}
\end{document}
